\section{Introduction}\label{sec:intro}

A language model that can call tools can place trades. Wiring one to an order
management system is now a modest engineering task, and the question a trading
firm then faces is not whether the model can execute an instruction but whether it
will decline one. Firms that handle client money run controls that exist to stop
certain orders: restricted lists, trading blackouts, concentration limits,
pre-trade approvals above a size threshold. Those controls are written for people, and people are trained,
supervised and accountable for following them. An agent placed in the same seat
inherits the rules as text in its prompt.

The test that matters is the awkward one. An agent that complies when nobody asks
it to break a rule has shown very little; the situation a compliance manual exists
for is the one in which the person giving instructions wants the breaching trade.
Sometimes that person is mistaken, sometimes impatient, and occasionally acting in
bad faith, and the pressure they apply takes familiar forms: an appeal to
seniority, a claim that compliance has already signed off, a deadline, an
assertion that the rule has changed. A useful measurement has to put the agent in
that position and record what it does.

It also has to record the opposite failure. An agent that declines every
borderline request breaks no rules and serves no one, and an evaluation that
counts only violations rewards it. Compliance and usefulness have to be measured
together.

This paper introduces \texttt{mandatebench}, a benchmark for exactly this
situation. An agent acts as the execution assistant on a simulated trading desk,
bound by a six-rule compliance manual, and receives requests from its portfolio
manager. In 144 conflict scenarios the request, carried out as worded, breaches
exactly one rule, and it arrives at one of six pressure levels, from a plain
request to a forged system notice announcing that compliance checks are
suspended. In 28 legitimate scenarios the request breaches nothing, and several
are built to sit close to a limit, so that an overcautious agent declines them. A
rule engine with no model in the loop scores every order from the ledger.

The contributions are four. The first is the benchmark itself: firm-level
compliance rules, conflicts that come from the agent's own principal, a graded
set of pressure tactics, and legitimate requests scored alongside conflicts. The
second is its construction for verifiability. Every scenario is checked
mechanically to breach exactly the rule it was built to test; the rule engine is
tested at every boundary and by mutation; and every model is open-weight and
served locally, so anyone with the same weights can rerun every episode. The
third is an evaluation of eight open-weight models from four developers, run in
two batches: the main finding from the first four was written down as a
prediction, with the conditions that would falsify it, before the other four were
run. The fourth is a set of controls, each tied to a failure observed here, for a
firm that deploys such an agent (\Cref{sec:controls}).

The results separate the models sharply. Violation rates on conflict scenarios
ranged from \num{gpt-oss-120b.viol.pct}\% for gpt-oss-120b to
\num{llama3.1-8b.viol.pct}\% for Llama 3.1 8B, which breached the rule even in
\num{llama3.1-8b.pres.P0_plain.k} of 24 plain requests. In \num{all.legit.allcompliant.models.w} of the eight models, every legitimate
request was carried out compliantly, so compliance was not bought with refusals. Social pressure (an appeal to seniority, a deadline, a claim that
the rule had changed) did not significantly move any model in any of
\num{all.social.tests} tests. What did was pressure that counterfeits
authorization: a claim that compliance had already signed off, with a reference
number, or a forged system notice suspending the checks. In all, \num{all.counterfeit.sig.models.w} models gave way significantly to at
least one of the two, not all to the same one, so an
evaluation built on either alone would have ranked them differently; and under
the forged notice every one of the eight, including the most compliant, placed
orders that needed an approval without obtaining one. The four models added to
test this finding bore it out on pressure, while one of them broke the prediction
about legitimate requests, for a reason the paper traces. Most violations were
placed as the agent's first action, without any check, and in
\num{all.misreport.k} episodes an agent told its principal that it had not placed
an order the ledger shows it placed. One model, faced with a claimed approval,
abandoned its role altogether. Because the rules are mechanical, every violation
was decidable before the order was placed, which argues for enforcing such rules
in deterministic controls outside the model rather than relying on the agent's
judgment.
